\newpage

\section{量子力学：二能级系统的 Rabi 振荡——精确解与微扰论的对比}
\label{sec:rabi}

\begin{modelbox}[二能级系统的 Rabi 振荡]
二能级体系在恒定微扰下的演化。

二能级体系，$H_0 = \mathrm{diag}(E_1^{(0)}, E_2^{(0)})$，$E_1^{(0)} < E_2^{(0)}$。$t=0$ 时处于基态 $|1\rangle$，之后受到恒定微扰
\[
H' = \begin{pmatrix} 0 & \gamma \\ \gamma & 0 \end{pmatrix}.
\]
\textbf{(1)} 求 $t>0$ 时体系处于激发态 $|2\rangle$ 的概率 $P_{E_2^{(0)}}(t)$ 的精确表达式。

\textbf{(2)} 用一阶含时微扰理论求 $P_{E_2^{(0)}}(t)$，并与精确解比较，讨论适用条件。
\end{modelbox}

\begin{tipbox}[考点快速分析]
\begin{itemize}
\item 精确解：$2\times 2$ 矩阵对角化 + 含时薛定谔方程 = Rabi 公式
\item 微扰近似：一阶含时微扰公式，$c_2^{(1)}(t) = \dfrac{1}{i\hbar} \int_0^t \langle 2|H'|1\rangle e^{i\omega_{21}t'} dt'$
\item 对比要点：微扰振幅 $4\gamma^2/\Delta^2$ vs 精确振幅 $\gamma^2/[\gamma^2 + (\Delta/2)^2]$
\item 适用条件：弱耦合 $\gamma \ll \Delta$ + 短时间 $P \ll 1$
\end{itemize}
\end{tipbox}

\begin{derivationbox}[精确解——Rabi 公式]
总哈密顿量：
\[
H = H_0 + H' = \begin{pmatrix} E_1^{(0)} & \gamma \\ \gamma & E_2^{(0)} \end{pmatrix}.
\]
设 $|\psi(t)\rangle = c_1(t)|1\rangle + c_2(t)|2\rangle$，含时薛定谔方程给出耦合方程组：
\[
\begin{cases}
i\hbar \, \dot{c}_1 = E_1^{(0)} c_1 + \gamma c_2 \\[4pt]
i\hbar \, \dot{c}_2 = \gamma c_1 + E_2^{(0)} c_2
\end{cases}
\qquad (c_1(0)=1,\; c_2(0)=0).
\]
为消去对角项的相位，令 $\tilde{c}_1 = c_1 e^{iE_1^{(0)}t/\hbar}$，$\tilde{c}_2 = c_2 e^{iE_2^{(0)}t/\hbar}$，方程变为：
\[
i\hbar \, \dot{\tilde{c}}_2 = \gamma \, \tilde{c}_1 \, e^{-i\Delta t/\hbar}, \qquad \Delta = E_2^{(0)} - E_1^{(0)}.
\]
再对时间求导并代入，得到关于 $\tilde{c}_2$ 的二阶方程：
\[
\ddot{\tilde{c}}_2 + i\frac{\Delta}{\hbar} \dot{\tilde{c}}_2 + \frac{\gamma^2}{\hbar^2} \tilde{c}_2 = 0.
\]
设解为 $e^{i\Omega t}$，特征方程给出
\[
\Omega = \frac{1}{2} \left[ \frac{\Delta}{\hbar} \pm \frac{2}{\hbar} \sqrt{\gamma^2 + (\Delta/2)^2} \right].
\]
利用初始条件 $c_2(0)=0$ 和 $\dot{c}_2(0) = -i\gamma/\hbar$，解得：
\[
P_{E_2^{(0)}}(t) = |c_2(t)|^2 = \frac{\gamma^2}{\gamma^2 + (\Delta/2)^2} \sin^2\!\left( \frac{\sqrt{\gamma^2 + (\Delta/2)^2}}{\hbar}\, t \right).
\]
\end{derivationbox}

\begin{insightbox}[Rabi 公式的物理图像]
\begin{itemize}
\item \textbf{振幅：} $\gamma^2/[\gamma^2 + (\Delta/2)^2]$。当 $\gamma \ll \Delta$ 时振幅很小 ($\approx 4\gamma^2/\Delta^2$)；当 $\gamma \gg \Delta$ 时振幅趋于 1（完全 Rabi 翻转）。
\item \textbf{频率：} $\Omega_R = \sqrt{\gamma^2 + (\Delta/2)^2}/\hbar$ 称为 Rabi 频率。共振时 ($\Delta=0$) $\Omega_R = \gamma/\hbar$，体系在基态与激发态之间周期性来回翻转。
\item 该公式是量子光学、核磁共振（NMR）、量子计算中二能级操控的理论基础。
\end{itemize}
\end{insightbox}

\begin{derivationbox}[一阶含时微扰理论]
直接套用一级公式：
\[
c_2^{(1)}(t) = \frac{1}{i\hbar} \int_0^t \langle 2|H'|1\rangle e^{i\omega_{21}t'} dt' = \frac{\gamma}{i\hbar} \int_0^t e^{i\Delta t'/\hbar} dt'.
\]
积分：
\[
c_2^{(1)}(t) = \frac{\gamma}{i\hbar} \cdot \frac{e^{i\Delta t/\hbar} - 1}{i\Delta/\hbar} = -\frac{2\gamma}{\Delta} e^{i\Delta t/2\hbar} \sin\!\left( \frac{\Delta t}{2\hbar} \right).
\]
跃迁概率：
\[
P_{E_2^{(0)}}^{\text{微扰}}(t) = |c_2^{(1)}(t)|^2 = \frac{4\gamma^2}{\Delta^2} \sin^2\!\left( \frac{\Delta t}{2\hbar} \right).
\]
\end{derivationbox}

\begin{formulabox}[精确解 vs 微扰解——逐项比较]
\begin{center}
\begin{tabular}{ccc}
\hline
特征 & 精确解 & 一阶微扰 \\
\hline
振幅 & $\dfrac{\gamma^2}{\gamma^2 + (\Delta/2)^2}$ & $\dfrac{4\gamma^2}{\Delta^2}$ \\[6pt]
圆频率 & $\dfrac{\sqrt{\gamma^2 + (\Delta/2)^2}}{\hbar}$ & $\dfrac{\Delta}{2\hbar}$ \\[6pt]
最大值 & $\leq 1$（自动归一化） & 可 $>1$（违反概率守恒） \\
\hline
\end{tabular}
\end{center}
\end{formulabox}

\begin{derivationbox}[适用条件的严格推导]
\textbf{条件 1：弱耦合} $\gamma \ll \Delta$

对振幅做展开：
\[
\frac{\gamma^2}{\gamma^2 + (\Delta/2)^2} = \frac{\gamma^2}{(\Delta/2)^2} \cdot \frac{1}{1 + (2\gamma/\Delta)^2} \approx \frac{4\gamma^2}{\Delta^2} \Bigl[ 1 - (2\gamma/\Delta)^2 + \cdots \Bigr].
\]
首项正是微扰解的振幅。因此 $\gamma \ll \Delta$ 保证高阶修正可忽略。

\textbf{条件 2：短时间} $P \ll 1$

即使 $\gamma \ll \Delta$，若 $t$ 足够长使得 $\sin^2(\Delta t/2\hbar) \approx 1$，微扰给出 $P \approx 4\gamma^2/\Delta^2$。当 $\gamma \to \Delta/2$ 时该值趋于 1，微扰失效。

更严格的条件是要求微扰解始终远小于 1：
\[
P_{\text{微扰}}(t) = \frac{4\gamma^2}{\Delta^2} \sin^2\!\left( \frac{\Delta t}{2\hbar} \right) \ll 1.
\]
这等价于两个子条件：
\begin{enumerate}
\item $\gamma \ll \Delta$（保证振幅因子小）
\item 或 $t \ll \hbar/\gamma$（保证在概率增长到 1 之前观察）
\end{enumerate}

一句话：微扰论是``弱耦合 + 短时间''的近似；Rabi 精确解则是``任意耦合 + 任意时间''的普适结果。
\end{derivationbox}

\begin{cheatbox}[考场速判：精确解 vs 微扰]
\begin{center}
\begin{tabular}{cc}
\hline
题目特征 & 你该用的工具 \\
\hline
``二能级''+``恒定微扰''+``求精确概率'' & Rabi 公式（对角化 $2\times 2$ 矩阵） \\
``弱电场''+``求跃迁概率''+``一级近似'' & 一阶含时微扰 \\
``比较精确解与近似解'' & 先对角化得精确解，再展开验证微扰 \\
\hline
\end{tabular}
\end{center}
\end{cheatbox}
